\subsection{Objetivo}
Probar una herramienta se ciberseguridad diseñada para atraer, engañar y observar el comportamiento de ciberdelincuentes con el fin observar como atacan y vulneran dispositivos en una red de computo, con la finalidad de es- tudiar los métodos utilizados.
\subsection{Qué es un Honeypot y T-Pot}
T-Pot es una plataforma con multiples honeypots montados en contenedores dockers, estos simulas diversos servicion vulnerables. (Del Hierro, 2025)\\
Con respecto a los T-Pop honeypot (Klein \& Pinkas, 2011) apoyan estas herramientas de ciberseguridad, como metodos de detección de vulnerabilidades al exponer una amplia variedad de sistemas simulados, con la finalidad de identificar como los ciberdelincuentes inspeccionan los sistemas, esto es crucial para reforzar la seguridad en la vida real para las personas y impresas.\\
Es importante mencionar que (Wang et al, 2023) está a favor de estas herramientas para aprender y comprender las tacticas que son empleadas por los ciberdelincuentes, proporcionando gran información sobre metodos, motivos y objetivos al atacar.
Otra de las ventajas de utilizar sistemas de simulación como T-Pot es que con los honeypots tienen la capacidad de proporcionar las herramientas que los atacastes usan en contra de las simulaciones esto es una gran ventaja, ya que se pueden actualizar las firmas de seguridad en los sistemas que se utilizan en la vida real. (Khan at al, 2023).\\
Como ventaja a poner tener las herramientas de los atacates, las empresas, organizaciones y gobiernos pueden hacer mejoras a sus sistemas, para adaptarse a las nuevas amenazas teniemdo como información base, la recopilación de datos que se pueden extraer de los honeypots (Symantec, 2023).
\subsection{Plataformas y Software}
La practica se realiso en el laboratorio de sistemas inteligentes y distribuidos en el edificio de tecnologias de la información en la universidad de la cañada.\\
Nodo Central (Honeypot): Computadora de escritorio Dell procesador Intel i7 12th generación, GPU Intel Iris, con 512GB de disco duro y 12GB de memoria RAM, con el sistema operativo Ubuntu Server 24.04.05 LTS.\\
\textbf{Nodo atacante 1:} Laptop HP Gaming Victus con procesador AMD Ryzen 7 8845HS, CPU Nvidia GeForce RTX 4050, con 1TB de almacenamiento de disco duro y 16GB de memoria RAM, con el sistema operativo Manjaro Gnome 25.\\
\textbf{Nodo atacante 2:} Laptop MSI GF63 Thin 11SC con procesador Intel Core i5-11400H, CPU Nvidia GeForce GTX 1650 Mobile Max-Q, con 1 TB de almacenamiento en disco duro y 256 GB en unidad SSD NVMe, además de 16 GB de memoria RAM, con el sistema operativo Manjaro Linux y entorno de escritorio KDE Plasma 6.5.\\
\textbf{Nodo atacante 3:} Laptop HP Gaming Victus con procesador AMD Ryzen 7 8845HS, CPU Nvidia GeForce RTX 4050, con 1TB de almacenamiento de disco duro y 16GB de memoria RAM, con el sistema operativo Manjaro Gnome 25.\\
\textbf{Nodo atacante 4:} Laptop HP Gaming Victus con procesador AMD Ryzen 7 8845HS, CPU Nvidia GeForce RTX 4050, con 1TB de almacenamiento de disco duro y 16GB de memoria RAM, con el sistema operativo Manjaro Gnome 25.
\subsection{Infraestructura de Hardware}
Docker \& Docker Compose: Se utilizó como el motor de virtualización a nivel de sistema operativo para desplegar la arquitectura de T-Pot. Esta tecnología permitió el aislamiento de los diversos honeypots en contenedores independientes, garantizando que las intrusiones capturadas no afectaran la estabilidad del sistema anfitrión.\\
Secure Shell (SSH): Protocolo de administración remota cifrada utilizado para la gestión del servidor. Durante la práctica, este servicio fue uno de los principales vectores de ataque analizados, registrando intentos de intrusión mediante técnicas de fuerza bruta y diccionarios de credenciales.\\
Htop: Herramienta de monitoreo interactivo de procesos en tiempo real. Se empleó para supervisar el impacto de los ataques de Denegación de Servicio (DoS) sobre los recursos de hardware, permitiendo observar picos de consumo en la CPU y memoria RAM durante las fases críticas de la actividad maliciosa.\\
Kibana: Interfaz de visualización de datos de código abierto integrada en el stack ELK. Actuó como el tablero de mando (dashboard) principal para el análisis forense, transformando los logs complejos de tráfico y ataques en gráficos, mapas de calor y cronologías visuales fáciles de interpretar.
\subsection{Infraestructura de Red}
Topología de Estrella Física: Conexión mediante un Switch Ethernet y cableado UTP Categoría 5e/6 para eliminar la inestabilidad de redes inalámbricas y garantizar la captura íntegra del tráfico.